% Figure: permitted daily exposure time versus A-weighted level, for the
% 3 dB (equal-energy) and 5 dB exchange rules, anchored at 8 h / 85 dB(A).
\begin{figure}[htbp]
\centering
\begin{tikzpicture}
\begin{axis}[
  width=12cm, height=7.5cm,
  xlabel={sound level (dB(A))}, ylabel={permitted time (h)},
  ymode=log,
  xmin=82, xmax=112, ymin=0.03, ymax=32,
  axis lines=left,
  domain=82:112, samples=200,
  legend pos=north east,
  legend cell align=left,
  every axis x label/.style={at={(current axis.right of origin)}, anchor=west},
  every axis y label/.style={at={(current axis.above origin)}, anchor=south},
  grid=major, grid style={enggray!20},
  tick label style={font=\scriptsize},
  label style={font=\small},
  legend style={font=\small},
]
% Equal-energy (3 dB) rule: T = 8 / 2^((L-85)/3).
\addplot[engblue, very thick] {8 / (2^((x-85)/3))};
\addlegendentry{$3$ dB rule (equal energy)}
% Older OSHA (5 dB) rule anchored at 8 h / 90 dB(A): T = 8 / 2^((L-90)/5).
\addplot[engred, very thick, dashed] {8 / (2^((x-90)/5))};
\addlegendentry{$5$ dB rule (older OSHA)}
% marker line at 8 hours
\addplot[enggray, thin, dashed] {8};
\node[enggray, font=\scriptsize, anchor=south west] at (axis cs:82.3,8.4)
  {full shift ($8$ h)};
\end{axis}
\end{tikzpicture}
\caption[Permitted exposure time under two exchange rules]{Permitted
daily noise exposure time as a function of A-weighted level, under the
$3$ dB equal-energy rule (blue, anchored at eight hours for
$85$ dB(A)) and the older $5$ dB OSHA rule (red, anchored at eight hours
for $90$ dB(A)). Each rule halves the permitted time for a fixed level
increment: three decibels under the energy rule, five under the OSHA
rule. The energy rule is the stricter one at high levels, because a
short burst of loud noise carries the same energy dose whether it is
scored on three or five decibels, and the two rules diverge sharply
above $95$ dB(A). Where a worker's level varies through the shift, the
total dose is the sum of partial doses $\sum t_i / T_i$, and a single
loud task can dominate the day even if it lasts minutes. The chart is
the exposure-accounting tool behind every hearing-conservation program:
it converts a level into a time budget rather than a pass-or-fail
threshold.}
\label{fig:vol03ch07:noise-dose}
\end{figure}
